Neural networks usually expose stochasticity at the final prediction, where a label or token is sampled from a softmax distribution. We ask what changes when sampling is moved into an intermediate representation instead. We compare deterministic inference, final-softmax sampling, categorical hidden activation sampling, stochastic activation pruning, Gaussian hidden noise, and Monte Carlo dropout in one controlled \cnn pipeline on \mnist and \fashionmnist. Across three seeds, stochastic methods are evaluated with $S \in \{1,5,20\}$ Monte Carlo samples and measured with accuracy, negative log likelihood, expected calibration error, predictive entropy, and sample variance. The main result is clear: hidden activation sampling is not equivalent to output sampling. A naive categorical sampler over penultimate hidden dimensions reduces accuracy by 19.13 percentage points on \mnist and 22.46 points on \fashionmnist at $S=20$, while increasing \ece from 0.002 to 0.254 on \mnist and from 0.010 to 0.202 on \fashionmnist. In contrast, stochastic activation pruning and Gaussian hidden noise preserve accuracy within 0.04 points on \mnist and within 0.15 points on \fashionmnist, while mostly increasing entropy and predictive variance. Final-softmax sampling preserves accuracy at $S=20$ but worsens empirical sampled-probability \nll. These results show that intermediate stochasticity is a representation-level intervention, not output sampling moved earlier, and that useful hidden sampling likely requires a better distribution and sample-aware training.
